\section{Cluster I: Substrate \& Emergence (Layers 1--4)}
\label{sec:cluster1}

The first cluster establishes the foundational observables and the initial dynamics by which structure and function emerge from elementary interactions. These four layers instantiate Axioms~\ref{ax:irred}--\ref{ax:emergence} and provide the substrate upon which all higher cognitive capacities are built.

%==============================================
\subsection{Layer 1: Foundational Observables (The Epistemic Singularity)}
\label{subsec:layer1}

The first layer of the \Apeiron{} Framework is the most fundamental: it defines the irreducible substrate from which all higher-level structure emerges. Without a principled notion of minimal granularity, any representational system succumbs to infinite regress or arbitrary stipulation. Layer~1 answers the question: \emph{What are the atoms of observation, and how can we know that they are indeed irreducible?}

We approach this question through the lens of \emph{multi-axial atomicity}. An observable is certified as atomic only if it is simultaneously recognized as such by four independent mathematical frameworks: Boolean algebra, measure theory, category theory, and information theory. Convergence across these frameworks provides robustness against the blind spots of any single formalism. The \texttt{MetaSpecification} associated with each observable encodes the weighting of these frameworks and a set of axioms that define the ontological grammar of the layer.

In the following, we present the formal definition of an \textit{Ultimate Observable}, the decomposition operators that test atomicity, the sublayers that provide supplementary perspectives, and the reference implementation that makes these concepts computationally executable.

\begin{figure}[htbp]
    \centering
    \adjustbox{width=\textwidth,center}{
        \begin{tikzpicture}[
            framework/.style={rectangle, draw=black, fill=green!10, minimum width=2.5cm, minimum height=0.8cm, align=center},
            operator/.style={rectangle, draw=black, fill=yellow!10, minimum width=3.2cm, minimum height=0.8cm, align=center, font=\footnotesize},
            arrow/.style={-{Stealth}, thick},
        ]
        \node[framework, minimum width=3cm, fill=blue!15] (obs) at (0,0) {Ultimate Observable\\$o \in \Obs$};
        
        \node[framework] (bool) at (-5,2.5) {Boolean\\Algebra};
        \node[framework] (meas) at (0,3.5) {Measure\\Theory};
        \node[framework] (cat) at (5,2.5) {Category\\Theory};
        \node[framework] (info) at (0,1) {Information\\Theory};
        
        \node[operator] (op_bool) at (-5,5) {BooleanDecompositionOperator\\$a \neq 0 \land \nexists b: 0<b<a$};
        \node[operator] (op_meas) at (0,6) {MeasureDecompositionOperator\\$\mu(A)>0 \land \forall B\subseteq A: \mu(B)\in\{0,\mu(A)\}$};
        \node[operator] (op_cat) at (5,5) {CategoricalDecompositionOperator\\Zero object (initial + terminal)};
        \node[operator] (op_info) at (0,4) {InformationDecompositionOperator\\$K(o) \approx |o|$ (incompressible)};
        
        \draw[arrow] (bool.north) -- (op_bool.south);
        \draw[arrow] (meas.north) -- (op_meas.south);
        \draw[arrow] (cat.north) -- (op_cat.south);
        \draw[arrow] (info.north) -- (op_info.south);
        
        \draw[arrow] (obs) -- (bool);
        \draw[arrow] (obs) -- (meas);
        \draw[arrow] (obs) -- (cat);
        \draw[arrow] (obs) -- (info);
        
        \node[framework, fill=orange!10] at (0,-2) {Atomicity Scores\\$\atomicity_{\text{bool}}, \atomicity_{\text{meas}}, \atomicity_{\text{cat}}, \atomicity_{\text{info}} \in [0,1]$};
        \draw[arrow] (bool.south) -- ++(0,-0.8) -| (-2,-2);
        \draw[arrow] (meas.south) -- ++(0,-0.5) -| (0,-2);
        \draw[arrow] (cat.south) -- ++(0,-0.8) -| (2,-2);
        \draw[arrow] (info.south) -- (0,-2);
        
        \node[framework, fill=purple!10, minimum width=4cm] at (0,-4) {MetaSpecification\\Weights, Axioms, Binary/Continuous};
        \draw[arrow] (0,-2.8) -- (0,-3.2);
        
        \end{tikzpicture}
    }
    \caption{Multi-axial atomicity validation in Layer~1. An \textit{Ultimate Observable} is tested against four independent mathematical frameworks.}
    \label{fig:decomposition}
\end{figure}

\textbf{Theoretical grounding.} Layer~1 is the irreducible substrate of representation. It operationalizes the philosophical notion of the \textit{Epistemic Singularity}---the point beyond which further decomposition yields no gain in discriminative or predictive content. This layer is anchored in four independent mathematical frameworks, each providing a necessary but individually insufficient condition for atomicity. Convergence across frameworks is mandated by Axiom~\ref{ax:multi}. The Boolean framework (Proposition~\ref{prop:bool-atom}) captures order-theoretic minimality; the measure-theoretic framework (Proposition~\ref{prop:meas-atom}) captures indivisibility with respect to a measure; the categorical framework (Proposition~\ref{prop:cat-atom}) captures universality through zero objects; and the information-theoretic framework (Proposition~\ref{prop:info-atom}) captures incompressibility as a proxy for Kolmogorov simplicity.

\paragraph{Epistemic Singularity as a Foundational Principle.}
The concept of an \emph{epistemic singularity} is drawn from the limits of knowledge representation: any observer-dependent granularity below this threshold carries no additional information that can influence the dynamics of the system. Formally, for any two putative sub‑observables $o', o''$ derived from $o$, the mutual information $I(o'; o'') \approx H(o)$ and the conditional entropy $H(o \mid o', o'') \approx 0$, indicating that the decomposition introduces redundancy without informational gain. This property is rigorously enforced by requiring that the \emph{integrated information} $\Phi(o) = 0$ for any proper partition of $o$, where $\Phi$ is computed using the formalism of Integrated Information Theory adapted to discrete observables.

\paragraph{Multi‑Axial Justification.}
Each framework tests a distinct aspect of indivisibility:
\begin{itemize}
    \item \textbf{Boolean Atomicity:} In a Boolean algebra $\mathcal{B}$, an element $a \neq 0$ is an atom if $0 < b \leq a$ implies $b = a$. For discrete values, the algebra is constructed from the set of all possible decompositions (e.g., prime factorization for integers, power set for sets). The operator $\decomp_{\text{bool}}$ returns the set of maximal proper divisors or subsets. Atomicity holds iff this set is empty.
    \item \textbf{Measure‑Theoretic Atomicity:} Given a measure space $(\Omega, \Sigma, \mu)$, a set $A \in \Sigma$ with $\mu(A) > 0$ is an atom if for every $B \subseteq A$ with $B \in \Sigma$, either $\mu(B)=0$ or $\mu(B)=\mu(A)$. The implementation approximates this by sampling subsets and checking measure invariance.
    \item \textbf{Categorical Atomicity:} In a category $\mathcal{C}$ with a zero object $0$, an object $X$ is \emph{simple} if it is isomorphic to a zero object. For $\mathbf{Set}_*$, the only zero object is the singleton set. The \texttt{CategoricalDecompositionOperator} attempts to factor $o$ as a product/coproduct of non‑trivial objects; if no such factorization exists, $o$ is categorically atomic.
    \item \textbf{Information‑Theoretic Atomicity:} An observable $o$ is atomic if its Kolmogorov complexity $K(o)$ satisfies $K(o) \geq |o| - c$, i.e., it is algorithmically random and therefore incompressible. In practice, zlib compression ratio serves as a computable upper bound on $K(o)$. For short strings, a conservative heuristic is applied.
\end{itemize}

\textbf{Formal definition.} Let $\Obs$ be the set of \textit{Ultimate Observables}. Each observable $o \in \Obs$ is a tuple $(\text{id}, \text{value}, \text{type}, \meta)$, where $\meta$ is a \texttt{MetaSpecification} instance. A decomposition operator is a map $\decomp_f: \Obs \to \mathcal{P}(\Obs)$ for a framework $f \in \{\text{bool}, \text{meas}, \text{cat}, \text{info}\}$. An observable $o$ is \emph{atomic} with respect to $f$ if $\decomp_f(o) = \{o\}$.

\begin{itemize}
    \item \textit{Boolean atomicity} ($f=\text{bool}$): $o$ is an atom in a Boolean algebra. For discrete values, the partial order is induced by divisibility (for integers) or set inclusion (for collections). Formally, $\decomp_{\text{bool}}(o) = \{o\}$ iff $\nexists b: 0 < b < o$.
    \item \textit{Measure atomicity} ($f=\text{meas}$): $o$ is an atom for a measure $\mu$. For the default counting measure, $\decomp_{\text{meas}}(o) = \{o\}$ iff $\mu(o) > 0$ and $\forall B \subseteq o: \mu(B) \in \{0, \mu(o)\}$. The implementation uses a configurable sampling limit to bound complexity.
    \item \textit{Categorical atomicity} ($f=\text{cat}$): $o$ is a zero object in its category. In the category $\mathbf{Set}_*$ of pointed sets, $\decomp_{\text{cat}}(o) = \{o\}$ iff $o$ is simultaneously initial and terminal. The \texttt{CategoricalDecompositionOperator} extracts subobjects and verifies the zero-object property.
    \item \textit{Information atomicity} ($f=\text{info}$): $o$ is incompressible. Using zlib as a proxy for Kolmogorov complexity, $\decomp_{\text{info}}(o) = \{o\}$ iff $|zlib(o)| \ge |o| - c$ for a small constant $c$. For ultra-short strings ($<10$ bytes) the proxy is unreliable; the implementation returns a conservative score of $1.0$ in such cases.
\end{itemize}

\textbf{Atomicity Score.} The overall atomicity score $\atomicity(o)$ is a weighted average of the framework-specific binary scores:
\[
\atomicity(o) = \sum_{f \in \{\text{bool}, \text{meas}, \text{cat}, \text{info}\}} \meta.\text{weights}[f] \cdot \mathbf{1}[\decomp_f(o) = \{o\}],
\]
where the weights are defined in the \texttt{MetaSpecification} and satisfy $\sum_f \meta.\text{weights}[f] = 1$. The \texttt{MetaSpecification} also encodes whether atomicity is interpreted as binary or continuous and maintains a set of axioms that govern the ontological grammar of the layer.

\paragraph{Observer Dependence and Contextual Atomicity.}
The \texttt{MetaSpecification} includes an \texttt{observer\_profile} that modulates the weights $w_f$ according to the observer's epistemic constraints. For instance, an observer with high measurement noise may down‑weight the measure‑theoretic criterion. Additionally, a \texttt{context} parameter allows the atomicity threshold to be relaxed in scenarios where approximate indivisibility suffices (e.g., coarse‑grained simulations). This implements Axiom~\ref{ax:observer} explicitly.

\textbf{Sublayers.} Although Layer~1 is conceptually indivisible, three supplementary perspectives are distinguished:
\begin{enumerate}[label=(\roman*)]
    \item \textit{Formal singularity}: The mathematical irreducibility captured by the decomposition operators.
    \item \textit{Empirical singularity}: The smallest unit detectable by a given measurement apparatus (e.g., Planck scale in physics, pixel in vision).
    \item \textit{Semantic singularity}: The minimal carrier of meaning in a symbolic system (e.g., phonemes, atomic propositions).
\end{enumerate}
Each sublayer corresponds to a distinct interpretation of the same underlying atomic structure, ensuring that the framework remains agnostic to the nature of the observer (human, machine, or hybrid).

\textbf{Implementation.} The reference implementation provides:
\begin{itemize}
    \item \texttt{UltimateObservable}: A dataclass encapsulating the observable's value, type, \texttt{MetaSpecification}, qualitative dimensions, relational embedding, and cached atomicity scores.
    \item \texttt{BooleanDecompositionOperator}, \texttt{MeasureDecompositionOperator}, \texttt{CategoricalDecompositionOperator}, \texttt{InformationDecompositionOperator}: Concrete implementations of the decomposition operators, registered in a global registry.
    \item \texttt{MetaSpecification}: A configuration object defining the primary principles, weights, axioms, and observer-dependence settings.
    \item \texttt{SelfProvingAtomicity}: A module that uses SymPy and Z3 to generate formal proofs of atomicity for critical observables.
    \item \texttt{AtomicityCache}: A thread‑safe cache with TTL and dirty‑flag invalidation to avoid redundant recomputation.
\end{itemize}

\paragraph{Algorithmic Verification of Atomicity.}
For observables that are critical to system integrity (e.g., those used in safety‑critical decisions), the \texttt{SelfProvingAtomicity} module attempts to construct a machine‑checkable proof of atomicity. Using Z3, it encodes the decomposition conditions as SMT formulas and attempts to prove unsatisfiability of the existence of a proper decomposition. For instance, for an integer $n$, the Boolean atomicity condition becomes:
\[
\forall d \in \mathbb{Z}^+,\; (d \mid n \land d \neq 1 \land d \neq n) \implies \text{False}.
\]
If the SMT solver returns \texttt{unsat}, the observable is certified as Boolean‑atomic. Similar encodings exist for the other frameworks.

\textbf{Computational Complexity.} The measure-theoretic test samples subsets to bound complexity at $\mathcal{O}(\text{max\_subsets})$, with exhaustive enumeration for $|o| \le 20$. The information-theoretic test runs in $\mathcal{O}(n)$ for input length $n$. Boolean and categorical tests are $\mathcal{O}(1)$ for typical scalar values. Caching of atomicity scores with dirty-flag invalidation avoids redundant recomputation.

\textbf{Axioms satisfied:} Axioms~\ref{ax:irred}, \ref{ax:multi}, \ref{ax:observer}. \\
\textbf{Research questions addressed:} RQ1, RQ2.

%==============================================
\subsection{Layer 2: Relational Emergence}
\label{subsec:layer2}

\paragraph{Introduction for the Reader.}
Having established the irreducible atoms of observation in Layer~1, we now confront a fundamental question: \emph{How does mere coexistence of atomic observables give rise to structure?} Layer~2 answers that structure is not an intrinsic property of things, but a pattern of relations. A single observable in isolation conveys no information beyond its own atomic value; it is only when observables co-occur, influence one another, or form higher-order assemblies that meaningful organization appears. This shift from \emph{substance} to \emph{relation} is the first step toward a non-anthropocentric ontology, because it avoids presupposing that objects exist prior to their interactions. In this section, we formalize the emergence of relational structure as a weighted hypergraph, prove fundamental properties of this representation, and connect it to research questions RQ2--RQ3 and Axioms~\ref{ax:relational} and~\ref{ax:emergence}.

\begin{figure}[htbp]
    \centering
    \adjustbox{width=0.8\textwidth,center}{
        \begin{tikzpicture}[
            node/.style={circle, draw=black, fill=blue!10, minimum size=0.8cm},
            edge/.style={-{Stealth}, thick, color=gray},
        ]
        \node[node] (a) at (0,0) {$o_1$};
        \node[node] (b) at (2,1) {$o_2$};
        \node[node] (c) at (4,0) {$o_3$};
        \node[node] (d) at (2,-1.5) {$o_4$};
        
        \draw[edge] (a) -- (b) node[midway, above left] {$0.8$};
        \draw[edge] (b) -- (c) node[midway, above right] {$0.6$};
        \draw[edge] (c) -- (d) node[midway, below right] {$0.9$};
        \draw[edge] (d) -- (a) node[midway, below left] {$0.3$};
        \draw[edge] (a) -- (c) node[midway, above] {$0.1$};
        \draw[edge] (b) -- (d) node[midway, right] {$0.4$};
        
        \node[draw, dashed, fit={(a) (b) (c) (d)}, inner sep=0.5cm, label=below:{\small Hypergraph $H = (\Obs, E)$}] {};
        
        \end{tikzpicture}
    }
    \caption{A relational hypergraph over a set of Layer~1 observables. Edge weights represent probabilistic relational strengths.}
    \label{fig:layer2}
\end{figure}

\subsubsection{Theoretical Grounding and Epistemological Shift}
Layer~2 operationalizes Axiom~\ref{ax:relational}: the identity of any higher-level entity is constituted by its relational position within a network of lower-level observables. This principle resonates with diverse intellectual traditions:
\begin{itemize}
    \item \textbf{Structuralism} (de Saussure, Lévi-Strauss): meaning arises from differences and relations within a system, not from intrinsic properties.
    \item \textbf{Network Science} \cite{barabasi2002}: complex systems are characterized by non-trivial topological features.
    \item \textbf{Bateson's Ecology of Mind} \cite{bateson1972}: information is a difference which makes a difference—a relation, not a thing.
    \item \textbf{Category Theory} \cite{maclane1971}: objects are defined up to isomorphism by their morphisms to and from other objects.
\end{itemize}
By adopting a purely relational substrate, we avoid importing human object-categories. The system does not see a ``chair''; it sees a set of atomic observables (pixels, edges, pressure values) whose co-occurrence patterns are statistically stable. The chair emerges as a cluster of relations, not as a labeled entity.

\subsubsection{Formal Definition: The Weighted Hypergraph of Relations}
\begin{definition}[Relational Structure]
Let $\Obs$ be the set of Ultimate Observables from Layer~1. A \emph{relational structure} is a triple $H = (\Obs, E, w)$ where:
\begin{itemize}
    \item $E \subseteq \bigcup_{k=1}^{K_{\max}} \Obs^k$ is a set of hyperedges. Each hyperedge $e = (o_{i_1}, \dots, o_{i_k}) \in E$ represents an irreducible $k$-ary relation among observables.
    \item $w: E \to [0,1]$ is a weight function assigning a probabilistic strength to each hyperedge, with the normalization $\sum_{e \in E} w(e) = 1$ (or more generally, $w(e)$ is the probability that $e$ is active in the current context).
\end{itemize}
\end{definition}

\begin{definition}[Induced Pairwise Adjacency]
For any two observables $o_i, o_j \in \Obs$, the \emph{effective adjacency} is defined as
\[
A_{ij} = \sum_{\substack{e \in E \\ \{o_i, o_j\} \subseteq e}} \frac{w(e)}{\binom{|e|}{2}} \cdot \mathbb{I}[\text{$e$ is active}],
\]
where the indicator reflects context-dependent activation. In the stationary limit, we may use $A_{ij} = \mathbb{E}[w(e) \mid o_i, o_j \in e]$. Note that we assume undirected hyperedges; the binomial coefficient $\binom{|e|}{2}$ counts the number of unordered pairs within the hyperedge.
\end{definition}

\begin{lemma}[Symmetry and Non-negativity]
$A_{ij} = A_{ji} \ge 0$ and $A_{ii} = 0$ (no self-loops).
\end{lemma}
\begin{proof}
Symmetry follows from the undirected nature of hyperedges; non-negativity from $w(e) \ge 0$. Self-loops are excluded because a hyperedge must contain at least two distinct observables to represent a relation.
\end{proof}

\begin{definition}[Relational Embedding]
A \emph{relational embedding} is a mapping $\mathbf{e}: \Obs \to \mathbb{R}^d$ such that the inner product $\langle \mathbf{e}_i, \mathbf{e}_j \rangle$ approximates $A_{ij}$ or a monotonic transformation thereof. The embeddings are obtained by minimizing a loss function, e.g.,
\[
\mathcal{L}_{\text{emb}} = \sum_{i,j} \left( A_{ij} - \sigma(\langle \mathbf{e}_i, \mathbf{e}_j \rangle) \right)^2 + \lambda \|\mathbf{e}_i\|^2,
\]
where $\sigma$ is the logistic sigmoid. This is a standard matrix factorization approach (see \cite{pennington2014glove}).
\end{definition}

\subsubsection{Probabilistic Semantics and Context-Dependence}
The weight $w(e)$ is not a fixed constant but a function of the system's state. Formally, we define a \emph{compatibility function} $\psi_e: \prod_{o \in e} \mathcal{V}_o \to \mathbb{R}$, where $\mathcal{V}_o$ is the value space of observable $o$. The weight is then computed via a softmax over all active hyperedges:
\[
w(e) = \frac{\exp(\psi_e(\mathbf{v}_{o_{i_1}}, \dots, \mathbf{v}_{o_{i_k}}))}{\sum_{e' \in E} \exp(\psi_{e'}(\dots))}.
\]
This allows the hypergraph to be \emph{dynamic}: the same set of observables can exhibit different relational strengths depending on contextual cues encoded in the compatibility functions. For instance, in a visual scene, the relation between two pixels might be strong only when a particular lighting condition is met.

\paragraph{Connection to Research Question RQ2.}
RQ2 asks: \emph{What is the minimal relational density required for the emergence of statistically significant functional motifs?} The above definition provides a precise measure of relational density: $\rho = \frac{1}{|\Obs|(|\Obs|-1)} \sum_{i \neq j} A_{ij}$. We will later prove a phase transition theorem (Theorem~\ref{thm:percolation}) linking $\rho$ to the appearance of giant connected components and motifs.

\subsubsection{Emergence of Structure: Weak Emergence and Excess Entropy}
\begin{definition}[Weak Emergence \cite{bedau1997}]
A property $P$ of the hypergraph $H$ is \emph{weakly emergent} iff:
\begin{enumerate}
    \item $P$ is nomologically determined by the micro-states (the observables and their pairwise relations).
    \item $P$ is irreducible to the micro-states in the sense that its description length is significantly smaller than the exhaustive listing of micro-states.
\end{enumerate}
\end{definition}
We quantify this via \emph{excess entropy} (also known as predictive information):
\[
E = I(\mathbf{X}_{\text{past}}; \mathbf{X}_{\text{future}}) - \sum_{i=1}^{|\Obs|} I(X_{i,\text{past}}; X_{i,\text{future}}),
\]
where $\mathbf{X}_t$ is the vector of observable states at time $t$. A positive excess entropy indicates that the whole predicts its future better than the sum of its parts, a signature of emergent organization.

\begin{proposition}[Excess Entropy as a Measure of Relational Emergence]
\label{prop:excess_entropy}
If $H$ is such that the mutual information between any two disjoint subsets of observables is strictly superadditive, then $E > 0$. Conversely, if all relations are pairwise and independent, $E = 0$.
\end{proposition}
\begin{proof}
By the chain rule for mutual information, $I(\mathbf{X}_{\text{past}}; \mathbf{X}_{\text{future}}) = \sum_i I(X_{i,\text{past}}; \mathbf{X}_{\text{future}} \mid \mathbf{X}_{<i,\text{past}})$. The excess $E$ isolates the synergistic components that cannot be attributed to individual variables.
\end{proof}

\subsubsection{Sublayers of Relational Structure}
We distinguish three sublayers, corresponding to increasing arity and complexity:
\begin{enumerate}
    \item \textbf{Dyadic Correlations} ($k=2$): Pairwise relations $r_{ij} = w(\{o_i, o_j\})$. Computed efficiently via co-occurrence counts or pointwise mutual information. This sublayer corresponds to the classical network representation.
    \item \textbf{Functional Dependencies}: Asymmetric, directed influences detected via Granger causality or transfer entropy. Formally, $o_i \rightarrow o_j$ if $I(o_j(t+1); o_i(t) \mid \Obs \setminus \{o_i, o_j\}) > \epsilon$. This sublayer anticipates the functional emergence of Layer~3.
    \item \textbf{Higher-order Motifs} ($k \ge 3$): Recurring sub-hypergraphs identified via tensor decomposition (e.g., non-negative Tucker decomposition) or frequent subgraph mining. These motifs are the seeds for functional clusters; their statistical overrepresentation is tested against a null model (e.g., configuration model for hypergraphs).
\end{enumerate}

\subsubsection{Mathematical Appendix: Spectral Properties and Percolation}
\begin{theorem}[Spectral Gap and Emergence of Clusters]
\label{thm:spectral}
Let $L = D - A$ be the unnormalized Laplacian of the effective adjacency matrix $A$, with $D_{ii} = \sum_j A_{ij}$. If the second smallest eigenvalue $\lambda_2(L) > 0$, the hypergraph is connected. Moreover, the number of connected components equals the multiplicity of the zero eigenvalue. As relational density $\rho$ increases, a phase transition occurs: for $\rho < \rho_c$, the graph is fragmented; for $\rho > \rho_c$, a giant component emerges.
\end{theorem}
\begin{proof}
Standard results from spectral graph theory \cite{chung1997spectral} apply to the weighted adjacency matrix. The critical density $\rho_c$ can be estimated via the Molloy-Reed criterion for random hypergraphs.
\end{proof}
This theorem directly addresses RQ2 by providing a mathematical condition for the emergence of macroscopic structure.

\paragraph{Online Hypergraph Maintenance.}
As new observables are added, the hypergraph is updated incrementally. The \texttt{DensityField} employs a streaming algorithm for hyperedge weight estimation, using exponential moving averages and reservoir sampling for motif discovery. This ensures that the relational structure remains responsive to new data without requiring full recomputation.

\textbf{Computational Complexity.} The pairwise influence computation is $\mathcal{O}(N^2)$ for $N$ observables. For large $N$, approximate nearest-neighbor methods (e.g., HNSW) reduce the effective cost to $\mathcal{O}(N \log N)$. Motif discovery is performed using sampling‑based methods with complexity $\mathcal{O}(|E| \cdot \text{sample\_size})$.

\textbf{Axioms satisfied:} Axioms~\ref{ax:relational}, \ref{ax:emergence}. \\
\textbf{Research questions addressed:} RQ2, RQ3.

%==============================================
\subsection{Layer 3: Functional Emergence}
\label{subsec:layer3}

\paragraph{Introduction for the Reader.}
The relational hypergraph of Layer~2 captures \emph{that} things co-occur and influence one another, but it does not yet tell us \emph{what they do}. Function is about reproducible transformation: a cluster of relations that, given certain inputs, reliably produces certain outputs. In this layer, we transition from structure to function, from correlation to causation. The key insight is that a function is an equivalence class of relational configurations—it can be realized by different underlying hypergraphs as long as the input-output mapping is preserved (multiple realizability). This abstraction is essential for the framework to recognize that a ``clock'' can be made of gears, quartz crystals, or atomic transitions; what matters is the functional role. Layer~3 fully instantiates Axiom~\ref{ax:emergence} and provides the operational answer to RQ3.

\begin{figure}[htbp]
    \centering
    \adjustbox{width=0.7\textwidth,center}{
        \begin{tikzpicture}[
            cluster/.style={ellipse, draw=black, fill=green!10, minimum width=2cm, minimum height=1cm, align=center},
            edge/.style={-{Stealth}, thick},
        ]
        \node[cluster] (f1) at (0,0) {$f_1$\\\footnotesize micro-function};
        \node[cluster] (f2) at (3,0) {$f_2$\\\footnotesize micro-function};
        \node[cluster, fill=orange!10] (meso) at (1.5,1.5) {$F_{\text{meso}}$\\\footnotesize meso-function};
        
        \draw[edge] (f1) -- (meso);
        \draw[edge] (f2) -- (meso);
        \draw[edge, dashed] (f1) to [bend left] (f2);
        \draw[edge, dashed] (f2) to [bend left] (f1);
        
        \node[draw, dashed, fit={(f1) (f2) (meso)}, inner sep=0.7cm, label=right:{\small Functional cluster}] {};
        
        \end{tikzpicture}
    }
    \caption{Functional emergence: clusters of hyperedges form stable functional entities, which can aggregate into higher-level functions.}
    \label{fig:layer3}
\end{figure}

\subsubsection{Theoretical Grounding: From Structuralism to Functionalism}
Layer~3 is grounded in the functionalist tradition in philosophy of mind \cite{putnam1967}, biology \cite{maturana1980}, and computer science \cite{turing1936}. The central tenet is that a function is defined by its causal role, not by its physical substrate. In the \Apeiron{} Framework, this translates to:
\begin{quote}
A \emph{functional entity} is a cluster of hyperedges $f \subseteq E$ such that the conditional probability distribution $P(O \mid I, f)$ is invariant under local perturbations of the hypergraph structure, and the cluster contributes non-redundantly to the system's overall predictive performance.
\end{quote}

\subsubsection{Formal Definitions: Functional Entities and Modularity}
\begin{definition}[Functional Entity]
Let $H = (\Obs, E, w)$ be a relational hypergraph. A \emph{functional entity} is a subset $f \subseteq E$ that satisfies:
\begin{enumerate}
    \item \textbf{Internal cohesion}: The modularity $Q(f) > \tau_{\text{mod}}$.
    \item \textbf{Causal autonomy}: Counterfactual ablation of $f$ significantly degrades a system-level objective $\mathcal{L}$ (see Definition~\ref{def:ablation}).
\end{enumerate}
\end{definition}

\begin{definition}[Hypergraph Modularity]
\label{def:modularity}
For a partition $\mathcal{P} = \{f_1, \dots, f_m\}$ of $E$, the modularity is
\[
Q(\mathcal{P}) = \frac{1}{2m} \sum_{f \in \mathcal{P}} \sum_{e,e' \in f} \left( w_e w_{e'} - \frac{\deg(e)\deg(e')}{2m} \right),
\]
where $m = \sum_{e \in E} w_e$, $\deg(e) = \sum_{e' \neq e} \text{overlap}(e,e') w_{e'}$, and $\text{overlap}(e,e')$ is the Jaccard index of the vertex sets. The modularity of a single cluster $f$ is its contribution to $Q(\mathcal{P})$.
\end{definition}

\begin{lemma}[Modularity and Community Detection]
Maximizing $Q$ over all partitions is NP-hard, but the Louvain algorithm provides a fast heuristic that scales to large hypergraphs. The resolution parameter $\gamma$ can be introduced to tune the granularity of detected functions: $Q_\gamma = Q - \gamma \sum_f |f|$.
\end{lemma}

\subsubsection{Counterfactual Ablation and Functional Significance}
\begin{definition}[Ablation Test]
\label{def:ablation}
Let $\mathcal{L}(H)$ be a system-level loss (e.g., prediction error on future observables). The \emph{ablation impact} of a candidate functional entity $f$ is
\[
\Delta_f \mathcal{L} = \mathcal{L}(H \setminus f) - \mathcal{L}(H),
\]
where $H \setminus f$ denotes the hypergraph with all edges in $f$ removed (or their weights set to zero). $f$ is \emph{functionally significant} if $\Delta_f \mathcal{L} > \delta$, where $\delta$ is a threshold determined by a null distribution (e.g., from random edge subsets of equal total weight).
\end{definition}

\begin{proposition}[Statistical Test for Functional Significance]
\label{prop:ablation}
Under the null hypothesis that $f$ is a random subset of edges with no functional role, the statistic $\Delta_f \mathcal{L}$ follows a distribution that can be estimated via bootstrap. A $p$-value $< \alpha$ rejects the null, confirming functional emergence.
\end{proposition}
\begin{proof}
By the central limit theorem for stationary processes, the bootstrap distribution of $\Delta_f \mathcal{L}$ converges to the true sampling distribution. We reject $H_0$ at level $\alpha$ if the observed $\Delta_f \mathcal{L}$ exceeds the $(1-\alpha)$ quantile of the bootstrap distribution.
\end{proof}

\subsubsection{Multiple Realizability and Equivalence Classes}
A crucial feature of functionalism is that the same function can be realized by different underlying structures. We formalize this as an equivalence relation.

\begin{definition}[Functional Equivalence]
Two hyperedge clusters $f, f' \subseteq E$ are \emph{functionally equivalent}, denoted $f \sim_F f'$, if for all input configurations $I$ (a set of observables external to the cluster), the output distributions satisfy
\[
D_{\text{JS}}\big( P(O \mid I, f), P(O \mid I, f') \big) \le \epsilon,
\]
where $D_{\text{JS}}$ is the Jensen-Shannon divergence. The \emph{functional type} $[f]$ is the equivalence class under $\sim_F$.
\end{definition}

This definition allows the system to recognize that a thermostat implemented via a bimetallic strip and one implemented via a digital sensor are instances of the same functional type. It also provides a foundation for the later emergence of ontological categories (Layer~9).

\subsubsection{Hierarchical Organization of Functions}
Functions are not flat; they organize into nested hierarchies. We define three canonical scales:
\begin{enumerate}
    \item \textbf{Micro-functions}: Small, dense clusters (typically 3-10 hyperedges) that act as computational primitives (e.g., AND-like gates, coincidence detectors). They are identified as strongly connected components in the directed influence graph.
    \item \textbf{Meso-functions}: Aggregations of micro-functions that exhibit meta-stability. They correspond to communities detected at a coarser resolution of modularity optimization. A meso-function may implement, for example, a regulatory feedback loop.
    \item \textbf{Macro-functions}: System-wide patterns that integrate multiple meso-functions. They define the global attractor landscape (Layer~4) and are the highest-level functional description before reaching the cognitive layers.
\end{enumerate}
The hierarchy is formalized by applying community detection recursively: the nodes of the level-$(l+1)$ graph are the functional entities of level $l$, and edge weights are the mutual information between their output streams.

\begin{proposition}[Hierarchical Consistency]
If $f$ is a meso-function composed of micro-functions $f_1, \dots, f_k$, then the ablation of $f$ produces a greater system-level impact than the sum of individual ablations of $f_i$ (synergy).
\end{proposition}
\begin{proof}
This follows from the non-additivity of the loss function: $\mathcal{L}(H \setminus f) - \mathcal{L}(H) > \sum_i (\mathcal{L}(H \setminus f_i) - \mathcal{L}(H))$, due to functional redundancy and compensation mechanisms within the meso-function.
\end{proof}

\subsubsection{Connection to Research Question RQ3}
RQ3 asks: \emph{How does the robustness of a functional entity scale with the modularity of its underlying relational cluster?} We answer this with a quantitative relationship.

\begin{theorem}[Modularity-Robustness Trade-off]
\label{thm:mod_robust}
For a functional entity $f$, let $\rho_f = \frac{1}{|f|} \sum_{e \in f} \text{robustness}(e)$ be the average edge robustness (measured by the variance of $w(e)$ under noise). Then there exists a positive correlation between $Q(f)$ and $\rho_f$, but beyond a critical modularity $Q^*$, robustness saturates and may even decline due to over-specialization (brittleness).
\end{theorem}
\begin{proof}
Empirically validated in Section~\ref{sec:validation}; theoretically, modularity concentrates information flow within $f$, reducing interference from outside noise, but extreme modularity isolates $f$ from corrective feedback.
\end{proof}

\subsubsection{Implementation Notes}
The reference implementation provides:
\begin{itemize}
    \item \texttt{FunctionalClusterer}: Applies Louvain with resolution tuning to the hypergraph.
    \item \texttt{AblationTester}: Computes $\Delta_f \mathcal{L}$ efficiently using influence propagation caching.
    \item \texttt{FunctionalEquivalence}: Maintains a hash of input-output signatures to group equivalent functions.
\end{itemize}

\textbf{Axioms satisfied:} Axiom~\ref{ax:emergence}, \ref{ax:relational}. \\
\textbf{Research questions addressed:} RQ3 directly, and provides the necessary substrate for RQ4.

%==============================================
\subsection{Layer 4: Dynamic Adaptation and Feedback}
\label{subsec:layer4}

\paragraph{Introduction for the Reader.}
The first three layers have built a static edifice: irreducible observables, their relational hypergraph, and functional clusters. But the world is not static; it flows. Layer~4 introduces \emph{time}—not as an external parameter, but as an emergent order from the causal dependencies among functional updates. Feedback loops close the loop: outputs of functions become inputs to other functions (or themselves), creating cycles of influence that enable self-regulation, learning, and the spontaneous emergence of complex dynamics. This layer fulfills Axiom~\ref{ax:time} and sets the stage for the adaptive intelligence of Cluster II. We will formalize the dynamics via stochastic differential equations, analyze stability and bifurcations, and directly address RQ4.

\begin{figure}[htbp]
    \centering
    \adjustbox{width=0.8\textwidth,center}{
        \begin{tikzpicture}[
            state/.style={circle, draw=black, fill=blue!10, minimum size=0.8cm},
            edge/.style={-{Stealth}, thick},
        ]
        \node[state] (s1) at (0,0) {$S_1$};
        \node[state] (s2) at (2,0) {$S_2$};
        \node[state] (s3) at (1,1.5) {$S_3$};
        
        \draw[edge] (s1) to [bend left] (s2);
        \draw[edge] (s2) to [bend left] (s1);
        \draw[edge] (s2) -- (s3);
        \draw[edge] (s3) -- (s1);
        
        \node at (1,-1) {\small $\frac{d\theta}{dt} = -\eta \nabla_\theta \mathcal{L}(\theta, x) + \sigma dW_t$};
        
        \end{tikzpicture}
    }
    \caption{Dynamic adaptation with feedback loops. State vectors evolve under stochastic differential equations, enabling self-regulation and non-linear dynamics.}
    \label{fig:layer4}
\end{figure}

\subsubsection{Theoretical Grounding: Time from Causality}
In the \Apeiron{} Framework, time is not a pre-existing Newtonian container. Instead, we adopt a causal set perspective: the order of events is defined by the partial order of functional influences. Let $G_{\text{inf}} = (\Func, E_{\text{inf}})$ be a directed graph where an edge $f_i \to f_j$ exists if $\frac{\partial S_j(t+1)}{\partial S_i(t)} \neq 0$ in expectation. A \emph{discrete time step} is defined as a synchronous update of all functions that have received new inputs since the last step. The \emph{temporal depth} $\tau(G)$ is the length of the longest directed path. Feedback loops create cycles, which we handle via simultaneous updates or by introducing delay differential equations.

\begin{definition}[State Space and Transition]
Each functional entity $f_k \in \Func$ is associated with a state vector $\theta_k \in \Theta_k \subseteq \mathbb{R}^{d_k}$. The joint state space is $\Theta = \bigoplus_k \Theta_k$. The dynamics are given by a stochastic transition function $\Phi: \Theta \times \mathcal{E} \to \Theta$, where $\mathcal{E}$ represents environmental inputs (other functions' outputs). For notational simplicity, we focus on a single function and drop the index $k$.
\end{definition}

\subsubsection{Stochastic Differential Equation for Adaptation}
We model the evolution of $\theta_t$ as an SDE derived from a potential $\mathcal{L}(\theta)$ (the objective functional):
\[
d\theta_t = -\eta \nabla_\theta \mathcal{L}(\theta_t) \, dt + \sigma \, dW_t,
\]
where:
\begin{itemize}
    \item $\eta > 0$ is the learning rate (adaptation speed).
    \item $\sigma \ge 0$ is the noise scale (exploration).
    \item $W_t$ is a standard Wiener process in $\mathbb{R}^d$.
\end{itemize}
This is the Langevin equation for the Gibbs distribution $p(\theta) \propto \exp(-\beta \mathcal{L}(\theta))$ with $\beta = 2\eta/\sigma^2$.

\paragraph{Choice of Objective $\mathcal{L}$.}
In the spirit of the Free Energy Principle \cite{friston2010}, we may set $\mathcal{L}(\theta) = \mathcal{F}(\theta) = D_{KL}[Q_\theta(s) \| P(s \mid o)] - \ln P(o)$, the variational free energy. Alternatively, for self-supervised learning, $\mathcal{L}$ could be a reconstruction error or contrastive loss. The specific form is less important than the existence of a gradient flow.

\begin{proposition}[Convergence to Stationary Distribution]
If $\mathcal{L}$ is $L$-smooth and strongly convex outside a compact set, the SDE admits a unique invariant measure $\pi(d\theta) = \frac{1}{Z} e^{-2\eta \mathcal{L}(\theta)/\sigma^2} d\theta$, and $\theta_t$ converges to $\pi$ exponentially fast in the Wasserstein-2 distance.
\end{proposition}
\begin{proof}
This is a standard result for overdamped Langevin dynamics \cite{roberts1996exponential}.
\end{proof}

\subsubsection{Feedback Loops: Homeostasis, Exploration, and Recursion}
We classify feedback according to its effect on the dynamics:

\begin{enumerate}
    \item \textbf{Homeostatic (Negative) Feedback}: Maintains internal variables within bounds. Linearized around an equilibrium $\theta^*$, the dynamics are $\frac{d\delta\theta}{dt} = \mathbf{J} \delta\theta$, where $\mathbf{J} = -\eta \nabla^2 \mathcal{L}(\theta^*)$. Stability requires all eigenvalues of $\mathbf{J}$ to have negative real parts.
    \item \textbf{Exploratory (Positive) Feedback}: Amplifies small fluctuations, helping escape local minima. This is implemented by making $\sigma$ state-dependent or by adding a repulsive drift term.
    \item \textbf{Recursive (Meta) Feedback}: The hyperparameters $\eta$ and $\sigma$ are themselves adapted based on the recent performance. For instance, $\frac{d\eta}{dt} = \gamma (\mathcal{L} - \mathcal{L}_{\text{target}})$, increasing the learning rate when loss is high.
\end{enumerate}

\begin{lemma}[Bifurcation Analysis]
Consider a simple one-dimensional system with feedback: $\frac{dx}{dt} = -x + \alpha \tanh(\beta x)$. For $\alpha \beta < 1$, the origin is a stable fixed point. At $\alpha \beta = 1$, a pitchfork bifurcation occurs, creating two stable fixed points. For larger $\alpha\beta$, the system becomes bistable. Adding noise induces stochastic resonance.
\end{lemma}
This lemma directly informs RQ4, which asks for the critical parameter regimes where feedback induces bifurcations.

\subsubsection{Lyapunov Exponents and Chaos}
For high-dimensional systems, the Lyapunov spectrum $\{\lambda_1 \ge \lambda_2 \ge \dots \ge \lambda_d\}$ characterizes the exponential divergence or convergence of nearby trajectories. A positive maximal Lyapunov exponent $\lambda_1 > 0$ is the signature of chaos.

\begin{definition}[Maximal Lyapunov Exponent]
\[
\lambda_1 = \lim_{t \to \infty} \lim_{\|\delta_0\| \to 0} \frac{1}{t} \ln \frac{\|\delta_t\|}{\|\delta_0\|},
\]
where $\delta_t$ is the evolution of an infinitesimal perturbation.
\end{definition}

\begin{theorem}[Chaos-Order Transition]
For the SDE with fixed $\eta, \sigma$, there exists a critical noise level $\sigma_c$ such that for $\sigma < \sigma_c$, the deterministic skeleton ($\sigma=0$) dominates and may exhibit chaos; for $\sigma > \sigma_c$, noise-induced order can stabilize irregular dynamics (stochastic resonance regime).
\end{theorem}
\begin{proof}
See \cite{arnold1998random} for random dynamical systems.
\end{proof}

\paragraph{Connection to RQ4.}
RQ4 asks: \emph{What are the critical parameter regimes at which a homeostatic feedback loop transitions into oscillatory or chaotic behavior?} Theorem~\ref{thm:bifurcation} provides a precise answer: the transition occurs when the leading eigenvalue of the Jacobian crosses the imaginary axis (Hopf bifurcation) or when the Lyapunov exponent becomes positive. These can be detected online by the \texttt{StabilityMonitor}.

\subsubsection{Implementation: Discretization and Event-Driven Simulation}
We discretize the SDE using Euler-Maruyama with step size $\Delta t$:
\[
\theta_{t+1} = \theta_t - \eta \nabla \mathcal{L}(\theta_t) \Delta t + \sigma \sqrt{\Delta t} \, \xi_t, \quad \xi_t \sim \mathcal{N}(0, I).
\]
To handle the large number of functions efficiently, we use an event-driven update scheme: a function $f$ is scheduled for update only when the accumulated influence from its neighbors exceeds a threshold $\epsilon_{\text{inf}}$, or when its internal prediction error surpasses $\epsilon_{\text{pred}}$. This mirrors the sparse communication in spiking neural networks.

\subsubsection{Mathematical Appendix: Stability Proofs}
\begin{lemma}[Lyapunov Function for Homeostasis]
For a system with negative feedback, $V(\theta) = \frac{1}{2}\|\theta - \theta^*\|^2$ is a Lyapunov function: $\frac{dV}{dt} \le -\alpha V$ for some $\alpha > 0$.
\end{lemma}
\begin{proof}
$\frac{dV}{dt} = (\theta - \theta^*) \cdot (-\eta \nabla \mathcal{L}(\theta)) \le -\eta \mu \|\theta - \theta^*\|^2$ by strong convexity of $\mathcal{L}$ near $\theta^*$.
\end{proof}

\begin{theorem}[Hopf Bifurcation in Feedback Systems]
\label{thm:bifurcation}
Consider a two-dimensional system $\frac{d\mathbf{x}}{dt} = \mathbf{F}(\mathbf{x}; \mu)$. If the Jacobian $D\mathbf{F}$ has a pair of complex conjugate eigenvalues $\alpha(\mu) \pm i\omega(\mu)$ with $\alpha(0)=0$, $\alpha'(0) > 0$, and $\omega(0) \neq 0$, then a Hopf bifurcation occurs at $\mu=0$, giving rise to a limit cycle.
\end{theorem}
This is a standard result in dynamical systems \cite{kuznetsov2004elements}. In our framework, $\mu$ could be the feedback gain $\alpha\beta$ from Lemma 1.

\subsubsection{Analogies and Examples}
\begin{itemize}
    \item \textbf{Thermostat}: A classic homeostatic feedback loop. The temperature sensor (functional entity) compares current temperature to setpoint and activates heater/AC. The SDE models the stochastic fluctuations due to sensor noise.
    \item \textbf{Population Dynamics (Lotka-Volterra)}: Predator-prey oscillations are an example of a limit cycle emerging from negative feedback (more prey $\to$ more predators $\to$ fewer prey $\to$ fewer predators $\to$ more prey). The framework can rediscover such cycles from raw observables.
    \item \textbf{Neural Avalanches}: In cortical networks, a balance of excitation and inhibition maintains the system near a critical point, producing power-law distributed avalanches. The SDE with state-dependent noise can replicate this criticality.
\end{itemize}

\textbf{Axioms satisfied:} Axiom~\ref{ax:time}, \ref{ax:emergence}. \\
\textbf{Research questions addressed:} RQ4 directly, and provides the dynamical foundation for RQ5--RQ7 in Cluster II.

\endinput